\section*{Capítulo \ref{ch:b2:biomolecules} --- Biomoléculas: aminoácidos, azúcares, lípidos y nucleótidos}

\begin{solution}{exo:b2:biomolecules:1}
\begin{itemize}
  \item Entre pH 3 y 9, el \omterm{def:b2:biomolecules:amino-acid}{ion dipolar} \ce{H3N+-CH(CH3)-COO-}.
  \item A pH 1, el catión \ce{H3N+-CH(CH3)-COOH}.
  \item A pH 12, el anión \ce{H2N-CH(CH3)-COO-}.
\end{itemize}
\end{solution}

\begin{solution}{exo:b2:biomolecules:2}
Glicina: $(2.35 + 9.78)/2 \approx 6.1$. Alanina: $(2.35 + 9.87)/2 \approx 6.1$.
% ledger: pka:glycine1, pka:glycine2, pka:alanine1, pka:alanine2
\end{solution}

\begin{solution}{exo:b2:biomolecules:3}
\ce{H2N-CH(CH3)-CO-NH-CH2-COOH}: la alanina aporta el extremo N-terminal, y la glicina, el C-terminal. Gly--Ala,
\ce{H2N-CH2-CO-NH-CH(CH3)-COOH}, es un isómero de constitución: el grupo amino libre está en la glicina.
\end{solution}

\begin{solution}{exo:b2:biomolecules:4}
\ce{NH2} a la izquierda: L. Prioridades \ce{NH2} > \ce{COOH} > \ce{CH2OH} > \ce{H}: la L-serina es
(\textit{S}).
\end{solution}

\begin{solution}{exo:b2:biomolecules:5}
Grupos: amonio N-terminal (unos 9), \omterm{def:b2:biomolecules:amino-acid}{cadena lateral} de la lisina (10.7), \omterm{def:b2:biomolecules:amino-acid}{cadena lateral} del aspartato (3.9),
grupo carboxílico C-terminal (unos 2). pH 1: $+1 +1 + 0 + 0 = +2$. pH 7: $+1 +1 -1 -1 = 0$. pH 12:
$0 + 0 - 1 - 1 = -2$ (la \omterm{def:b2:biomolecules:amino-acid}{cadena lateral} de la lisina, 1.3 unidades por debajo de pH 12, es mayoritariamente neutra).
\end{solution}

\begin{solution}{exo:b2:biomolecules:6}
El cloruro de tionilo y el calor atacarían también otros grupos y, sobre todo, el \omterm{def:b2:acyl-substitution:derivative}{cloruro de acilo}, muy reactivo,
de un aminoácido pierde fácilmente su configuración (el hidrógeno $\alpha$ se vuelve ácido): el \omterm{def:b2:biomolecules:peptide}{péptido} quedaría
parcialmente racemizado. Un \omterm{def:b2:biomolecules:peptide}{agente de acoplamiento} activa el ácido a temperatura ambiente, en condiciones suaves.
\end{solution}

\begin{solution}{exo:b2:biomolecules:7}
O5 unido a C2: un anillo de C2, C3, C4, C5 y O, de cinco átomos (una furanosa). En \omterm{def:b2:biomolecules:anomer}{proyección de Haworth}, el O del anillo
atrás, C2 a la derecha con el \ce{CH2OH} (C1) hacia abajo y el \ce{OH} hacia arriba ($\beta$); el \ce{OH} de C3 hacia arriba (a la izquierda en
Fischer), el \ce{OH} de C4 hacia abajo, y C5 con el \ce{CH2OH} (C6) hacia arriba.
\end{solution}

\begin{solution}{exo:b2:biomolecules:8}
$x_\alpha = (80.2 - 52.8)/(150.7 - 52.8) = 27.4/97.9 \approx 0.28$: 28~\% de $\alpha$, 72~\% de $\beta$.
\end{solution}

\begin{solution}{exo:b2:biomolecules:9}
La maltosa conserva un hemiacetal libre en su segunda glucosa, que se abre a un aldehído: es reductora. En la sacarosa
ambos \omterm{def:b2:biomolecules:anomer}{carbonos anoméricos} están comprometidos en el \omterm{def:b2:biomolecules:glycoside}{enlace glicosídico}: sin hemiacetal, no es reductora. El ácido diluido
hidroliza el \omterm{def:b2:biomolecules:glycoside}{enlace glicosídico}, un acetal: se liberan glucosa y fructosa, y ambas son reductoras.
\end{solution}

\begin{solution}{exo:b2:biomolecules:10}
La trioleína, \qty{885.45}{g/mol}, consume tres \ce{KOH} (\qty{56.105}{g/mol}): $3 \times 56.105/885.45 =
0.190$~g por g, un índice de \omterm{def:b2:acyl-substitution:saponification}{saponificación} de 190. Cadenas más cortas implican una masa molar menor para los mismos
tres grupos éster: más moles de grasa por gramo, más KOH.
% ledger: aw:C, aw:H, aw:O, aw:K
\end{solution}

\begin{solution}{exo:b2:biomolecules:11}
Protéjase la amina de la glicina como carbamato de \textit{terc}-butoxicarbonilo, y el ácido de la alanina como
éster metílico. Acóplese con una carbodiimida. Elimínense el carbamato con ácido y el éster con base diluida
(condiciones ortogonales, de modo que cualquiera de los dos extremos puede liberarse por separado): Gly--Ala.
\end{solution}

\begin{solution}{exo:b2:biomolecules:12}
5$'$-ACGCAT-3$'$ (leída en antiparalelo: A se aparea con T, y G con C). Pares: A--T, T--A, G--C, C--G, G--C, T--A:
$3 \times 2 + 3 \times 3 = 15$ enlaces de hidrógeno.
\end{solution}

\begin{solution}{pb:b2:biomolecules:1}
\textbf{1.} Ácido L-aspártico, L-fenilalanina y metanol.
\textbf{2.} En ambos, \ce{COOH} arriba, la \omterm{def:b2:biomolecules:amino-acid}{cadena lateral} abajo (\ce{CH2COOH}, \ce{CH2C6H5}), y el \ce{NH2}
a la izquierda.
\textbf{3.} (\textit{S}) para ambos.
\textbf{4.} Dos estereocentros: cuatro estereoisómeros.
\textbf{5.} \ce{H2N-CH(CH2COOH)-CO-NH-CH(CH2C6H5)-COOCH3}: el \omterm{def:b2:biomolecules:peptide}{enlace peptídico} une los dos
restos; el éster es el \ce{COOCH3} del extremo C-terminal.
\textbf{6.} Los receptores del gusto son quirales: se unen a un estereoisómero y no a su imagen especular ni a sus
diastereoisómeros, como decía la introducción.
\textbf{7.} El \ce{NH3+} N-terminal y el \ce{COOH} de la \omterm{def:b2:biomolecules:amino-acid}{cadena lateral} del ácido aspártico; el ácido C-terminal
es un éster metílico, sin hidrógeno ácido.
\textbf{8.} A pH 2: \ce{NH3+} ($+1$), \ce{COOH} de la \omterm{def:b2:biomolecules:amino-acid}{cadena lateral} (0): $+1$.
\textbf{9.} A pH 7: $+1 - 1 = 0$.
\textbf{10.} A pH 12: $0 - 1 = -1$.
\textbf{11.} Entre los dos $\mathrm pK_a$ de los grupos que cambian: $(3.9 + 9.9)/2 \approx 6.9$.
\textbf{12.} En el dipéptido el carbono de la amina lleva una amida en lugar de un carboxilato: los
grupos y cargas vecinos cambian, y con ellos los $\mathrm pK_a$ (el amonio N-terminal de un
\omterm{def:b2:biomolecules:peptide}{péptido} es claramente más ácido que el de un aminoácido libre). El modelo da la forma, no los
valores exactos.
\textbf{13.} El ácido y la amina forman primero una sal; calentar daría mezclas (Asp--Asp, cadenas más largas,
el grupo carboxílico equivocado) y racemizaría los estereocentros.
\textbf{14.} Su grupo amino y el grupo carboxílico de su \omterm{def:b2:biomolecules:amino-acid}{cadena lateral}.
\textbf{15.} Convierte el \ce{OH} del ácido en un buen grupo saliente, un derivado activado que la
amina ataca a temperatura ambiente (sustitución de acilo).
\textbf{16.} El $\beta$-dipéptido, unido a través de la \omterm{def:b2:biomolecules:amino-acid}{cadena lateral}: un isómero del aspartamo, no el aspartamo.
\textbf{17.} La amina como carbamato de benciloxicarbonilo y la \omterm{def:b2:biomolecules:amino-acid}{cadena lateral} como éster bencílico: ambos se
eliminan a la vez por hidrogenólisis sobre paladio, que deja intacto el éster metílico (ortogonal
a él).
\textbf{18.} Un ácido activado tiene un hidrógeno $\alpha$ ácido; la base lo arranca y el estereocentro
se racemiza.
\textbf{19.} Protéjase el ácido aspártico (carbamato sobre N, éster bencílico en la \omterm{def:b2:biomolecules:amino-acid}{cadena lateral}); actívese el
$\alpha$-\ce{COOH} libre con el \omterm{def:b2:biomolecules:peptide}{agente de acoplamiento}; añádase el éster metílico de la L-fenilalanina; hidrogenolícese.
\textbf{20.} El éster metílico y el \omterm{def:b2:biomolecules:peptide}{enlace peptídico}; primero el éster, pues una amida es mucho menos reactiva.
\textbf{21.} El dipéptido Asp--Phe (ácido libre) y metanol.
\textbf{22.} La molécula dulce se consume: sus productos de hidrólisis no la sustituyen.
\textbf{23.} $14 \times 12.011 + 18 \times 1.008 + 2 \times 14.007 + 5 \times 15.999 =
\qty{294.307}{g/mol}$.
\textbf{24.} $0.180/294.307 = \qty{6.116e-4}{mol}$, \qty{0.612}{mmol}.
\textbf{25.} Un metanol por aspartamo: $0.6116 \times 32.042 = \boldsymbol{\approx \qty{19.6}{mg}}$ de
metanol.
methanol.
\end{solution}
